\section{AAA 范式：把评测器本身做成 Agent}

本讲提出的 `Identified Agent Assessments (AAA)` 是一个关键思路：不再把 benchmark 写成静态 harness，而是把 ``评测器'' 本身做成 assessor agent。assessor 通过标准协议暴露工具接口，被测 agent 只需遵循统一接口（例如 MCP 风格）即可接入，从而降低适配复杂度。

\begin{importantbox}{AAA 的工程收益}
\begin{itemize}
    \item 从 ``每个 benchmark 写一套胶水'' 变成 ``一次接入，复用多评测器''；
    \item 更容易做 reproducible 的对抗评测；
    \item 更容易扩展 arena 模式（多 agent 竞争/对抗）；
    \item 更容易做持续监控和回归测试。
\end{itemize}
\end{importantbox}

课程还提到 Agent Beats 这类开放平台，目标是让社区持续提交新环境、新 benchmark、新 agent，并在统一框架下比较能力与风险。这种平台化路线对于安全研究尤为重要，因为很多漏洞只有在开放生态和持续对抗中才会暴露。

\begin{warningbox}{评测平台本身也会成为攻击目标}
一旦评测影响资源分配和声誉，平台就会面临 benchmark gaming、test leakage 和评测规避策略。因此评测平台需要独立的安全治理机制。
\end{warningbox}

\subsection{本章小结}
AAA 的核心价值是把评测标准化问题前置解决，降低集成摩擦并提升可复现性。它是 ``安全能力工程化'' 的基础工具，而不是附属组件。

